% Source-ready Evaluation replacement for the current paper source project.
% The main LaTeX source is absent from this repository; see paper_update_blocker.md.
% All numerical claims below come from results/evaluation_v1/processed CSV files.

\section{Evaluation}
\label{sec:evaluation}

\subsection{Methodology}
We evaluate Static-Sym, Sym-OCS, and DRAC on ordered DP, PP, and DP+PP
communication graphs.  Sym-OCS receives the same endpoint channel budget,
ordered node sequence, and reconfiguration-delay model as DRAC, but enforces
$X_{uv}=X_{vu}$.  Unless otherwise stated, the deterministic simulator uses a
25-Gb/s fixed path, 100 Gb/s per OCS connection unit, seed 7, a 0.5-ms
reconfiguration delay, and realization tolerance $\epsilon=0.1$.  Simulator
outputs are not hardware measurements.

\subsection{Accuracy of Ordered-Demand Profiling}
\label{sec:eval-profiler}
\textbf{Measurement TODO.} The artifact does not contain raw directional NIC
counters.  Figure~\ref{fig:profiler-accuracy} therefore reports the missing-input
state and does not present simulated traffic as measured traffic.  Once the raw
counters are supplied, the checked-in loader evaluates payload-only and
payload-plus-calibration predictions using weighted absolute, median, P95,
main-direction, reverse/control-direction, and dominant-direction errors.

\begin{figure}[t]
  \centering
  \includegraphics[width=\columnwidth]{results/evaluation_v1/figures/figure_6_profiler_accuracy.pdf}
  \caption{Ordered-demand profiler validation status.  Measured directional NIC
  counters are not present in this artifact; no simulated value is labeled as a
  measurement.}
  \label{fig:profiler-accuracy}
\end{figure}

\subsection{End-to-End Communication Performance}
\label{sec:eval-e2e}
Figure~\ref{fig:e2e} reports communication and reconfiguration time from the
revised simulator.  At eight channels per endpoint, DRAC reduces the DP time
from 2.842 ms for Sym-OCS to 1.464 ms.  The PP and mixed cases are negative
results: DRAC takes 7.452 ms and 8.585 ms, versus 5.952 ms and 6.632 ms for
Sym-OCS, because the selected schedules pay more reconfiguration overhead.
These results motivate further PP scheduling work and prevent an unsupported
no-harm claim.

\begin{figure*}[t]
  \centering
  \includegraphics[width=\textwidth]{results/evaluation_v1/figures/figure_7_end_to_end_performance.pdf}
  \caption{End-to-end communication time for DP, PP, and DP+PP ordered
  communication graphs.  All three schemes use the same graph, resource budget,
  and delay model.}
  \label{fig:e2e}
\end{figure*}

\subsection{Directional Target Generation and Segmentation}
\label{sec:eval-segmentation}
The validation compares the closed-form per-node target against a numerical
reference and exhaustive small instances.  Figure~\ref{fig:segmentation} shows
the segmentation trade-off.  At low delay the DP solution retains four
representative configurations for this trace; at 10 ms it collapses to one
configuration, matching the expected OneConfig limit.

\begin{figure*}[t]
  \centering
  \includegraphics[width=\textwidth]{results/evaluation_v1/figures/figure_8_segmentation.pdf}
  \caption{Segmentation total cost and selected segment count as the
  reconfiguration delay $\delta$ changes.}
  \label{fig:segmentation}
\end{figure*}

\subsection{Sparse Integer OCS Realization}
\label{sec:eval-realization}
Figure~\ref{fig:realization} compares FloorOnly, NearestRounding,
FillAllResidual, DRACSparse, and a small exhaustive oracle.  DRACSparse stops
when the slowdown tolerance is met and then reverse-prunes; it does not consume
all remaining ports merely to reduce target residual.

\begin{figure*}[t]
  \centering
  \includegraphics[width=\textwidth]{results/evaluation_v1/figures/figure_9_realization_tradeoff.pdf}
  \caption{Realized slowdown and schedule-wide stable channels versus
  realization tolerance $\epsilon$.}
  \label{fig:realization}
\end{figure*}

\subsection{Schedule-Wide Resource Compaction}
\label{sec:eval-compaction}
The planner reserves the maximum outgoing and incoming degree over the complete
schedule, rather than exposing ports that happen to be idle in one segment.
For the evaluated eight-channel configurations, the schedule-wide peak reserves
48 of 64 directional channels and stably exposes 16.

\begin{figure*}[t]
  \centering
  \includegraphics[width=\textwidth]{results/evaluation_v1/figures/figure_10_schedule_compaction.pdf}
  \caption{Schedule-wide stable reserved and exposed directional channels.}
  \label{fig:compaction}
\end{figure*}

\begin{figure}[t]
  \centering
  \includegraphics[width=\columnwidth]{results/evaluation_v1/figures/figure_11_iso_performance_pool.pdf}
  \caption{Minimum stable channel pool at the communication time of the
  four-channel Sym-OCS reference.  Independent Tx/Rx and bidirectional-bundle
  accounting are shown separately.}
  \label{fig:iso-pool}
\end{figure}

\subsection{Sensitivity and Offline Planning Overhead}
\label{sec:eval-overhead}
The artifact records graph parsing, ordered profiling, per-node target
generation, service-matrix construction, candidate segment costs, dynamic
programming, sparse realization, and schedule compaction separately.  Runtime
values are machine-dependent and are regenerated into
Table~\ref{tab:target-runtime}; no runtime is manually embedded here.

\begin{table}[t]
  \centering
  \caption{Target solver correctness and offline planning runtime.}
  \label{tab:target-runtime}
  \input{results/evaluation_v1/tables/target_solver_validation.tex}
\end{table}

\subsection{Summary}
The revised evaluation validates node-level target optimality, target-sequence
segmentation, tolerance-driven sparse realization, and schedule-wide resource
compaction.  It also identifies two limits of the current artifact: measured
profiler counters are missing, and PP-heavy schedules do not yet satisfy the
desired no-harm behavior.  We therefore do not reuse the legacy skew-threshold,
share-error, residual-gap, or headline speedup results.

% Abstract/Introduction/Conclusion integration note:
% Do not copy legacy headline numbers. State that profiling measurement is
% pending and that the current deterministic simulator demonstrates a DP benefit
% but exposes unresolved PP/mixed reconfiguration overhead. Integrate only after
% the complete paper source is restored and cross-references can be compiled.

